\section{Benchmark-Skript für datenbankbasierte API-Endpunkte (5123026)}
\label{sec:benchmark_script}

Zur Bewertung der Datenbank-Performance wurde das Skript \texttt{benchmark\_db\_endpoints.py} erstellt. Es misst ausschließlich lesende \texttt{GET}-Endpunkte, damit ein Benchmark-Lauf keine Daten erzeugt, verändert oder löscht. Dadurch bleiben die Messungen wiederholbar und können genutzt werden, um Performanceverbesserungen wie Indizes, Entity Graphs, Caching oder Pagination miteinander zu vergleichen. Der Fokus liegt nicht auf einem isolierten SQL-Test, sondern auf der tatsächlichen Laufzeit der Backend-Endpunkte inklusive HTTP-Schicht, Controller, Service, Repository, ORM und Datenbankzugriff.

Zu Beginn prüft das Skript, ob das Backend bereits läuft oder gestartet werden muss. Standardmäßig wird bei Bedarf die Spring-Boot-Anwendung über Gradle gestartet. Danach wartet das Skript, bis ein einfacher datenbankbasierter Endpunkt erreichbar ist. Diese Bereitschaftsprüfung ist wichtig, weil ein zu früh gestarteter Benchmark nicht die Performance der Datenbankzugriffe messen würde, sondern den noch laufenden Startprozess der Anwendung.

Vor den eigentlichen Messungen ermittelt das Skript repräsentative Beispieldaten. Dafür werden die Collection-Endpunkte der zentralen Tabellen abgefragt, etwa Departments, Stations, Rooms, Bookings, Patients, Doctors, Nurses, Drugs, Doses, Medications und Diagnoses. Aus den ersten Rückgabewerten werden benötigte IDs oder Fachwerte extrahiert, zum Beispiel eine Raum-ID, eine Patienten-ID, ein Stockwerk oder eine Stations-ID. Dieser Schritt ist notwendig, weil viele Endpunkte keine statischen Pfade besitzen. Ein Aufruf wie \texttt{/api/patients/\{id\}/diagnoses} kann nur sinnvoll gemessen werden, wenn vorher eine existierende Patienten-ID bekannt ist.

Anschließend werden die Endpunktvorlagen mit den gefundenen Beispieldaten gefüllt. Wenn für einen Endpunkt keine passenden Beispieldaten gefunden werden, wird er übersprungen und im Ergebnis dokumentiert. Dadurch bleibt der Benchmark robust gegen unvollständige Testdaten. Gleichzeitig wird sichtbar, ob ein Messergebnis fehlt, weil die Datenbasis keinen geeigneten Datensatz enthält.

Vor jeder Messreihe führt das Skript einen Warmup durch. Beim Warmup wird derselbe Endpunkt mehrfach aufgerufen, ohne diese Zeiten in die spätere Auswertung aufzunehmen. Dieser Schritt ist für aussagekräftige Messungen wichtig, weil die erste Anfrage oft Sonderkosten enthält: Die JVM optimiert Code erst zur Laufzeit, Hibernate initialisiert Metadaten und Query-Pläne, Verbindungspools müssen aufgebaut werden, und die Datenbank sowie das Betriebssystem füllen ihre Caches. Ohne Warmup würden die Messergebnisse diese Startkosten mit der eigentlichen stabilen Laufzeit vermischen. Der Warmup bringt die Anwendung in einen realistischeren Betriebszustand, in dem wiederholte Anfragen mit bereits initialisierten Strukturen verarbeitet werden.

Nach dem Warmup beginnt die eigentliche Messphase. Für jeden Endpunkt wird eine konfigurierbare Anzahl an Requests ausgeführt. Standardmäßig sind das 30 Messungen pro Endpunkt nach 5 Warmup-Aufrufen. Jede Messung erfasst die Dauer des gesamten Requests in Millisekunden, den HTTP-Status und die Größe der gelesenen Antwort. Erfolgreiche Antworten werden für die Laufzeitstatistik verwendet; Fehler oder Timeouts bleiben dennoch in den Rohdaten erhalten, damit instabile Endpunkte nicht unbemerkt aus der Auswertung verschwinden.

Die erhobenen Einzelwerte werden am Ende pro Endpunkt zusammengefasst. Neben dem Durchschnitt werden Median, Minimum, Maximum, Standardabweichung und das 95. Perzentil berechnet. Der Durchschnitt zeigt die allgemeine Tendenz, kann aber durch einzelne Ausreißer verzerrt werden. Der Median beschreibt die typische Antwortzeit stabiler. Das 95. Perzentil ist für Performance-Bewertungen besonders nützlich, weil es zeigt, wie langsam die langsamsten regulären Anfragen werden, ohne nur den absoluten Extremwert zu betrachten. Die Statusverteilung wird ebenfalls gespeichert, damit ein schneller Endpunkt nicht fälschlich als erfolgreich bewertet wird, wenn viele Requests nur Fehlerantworten liefern.

Die Ergebnisse werden in zwei Formaten ausgegeben. Die JSON-Datei enthält Metadaten, gefundene Beispieldaten, übersprungene Endpunkte und alle Einzelmessungen. Sie eignet sich für eine detaillierte technische Nachvollziehbarkeit. Die CSV-Datei enthält eine kompakte Zusammenfassung pro Endpunkt und kann einfacher in Tabellenkalkulationen oder Diagrammen verglichen werden. Für die Bewertung der Datenbankoptimierungen ist besonders relevant, dass mehrere Benchmark-Läufe mit denselben Parametern vor und nach einer Änderung verglichen werden können.

Das Skript verbindet die beschriebenen Performanceoptimierungen mit messbaren Effekten im Anwendungskontext. Ein Index ist erst dann fachlich überzeugend, wenn er die Endpunkte beschleunigt, die das System tatsächlich nutzt. Der Benchmark macht diese Wirkung sichtbar, ohne die Datenbank direkt mit isolierten SQL-Kommandos zu testen. Dadurch wird gemessen, was für Benutzer und andere Systemkomponenten relevant ist: die Antwortzeit der datenbankgestützten API.
